%==============================================================================
\section{Basic Concepts --- Các khái niệm cơ bản}
%==============================================================================

\begin{dongco}[title={Mục tiêu}]
Bài này giới thiệu các khái niệm cơ bản trong học máy: dữ liệu, đặc trưng, mô hình, huấn luyện/kiểm thử, overfitting/underfitting, và khi nào nên dùng học máy.
\end{dongco}

\subsection{Các thành phần cốt lõi}

\subsection{Overfitting}

\begin{dinhnghia}[title={Overfitting (Quá khớp)}]
Overfitting xảy ra khi mô hình quá phức tạp và khớp dữ liệu huấn luyện “quá sát”.
Hệ quả là mô hình hoạt động kém trên dữ liệu mới vì đã “chuyên biệt” cho tập train.
\end{dinhnghia}

\begin{phuongphap}[title={Cách giảm overfitting}]
\begin{itemize}
  \item Dùng tập validation để đánh giá trong quá trình phát triển.
  \item Dùng \textbf{regularization} để làm mô hình đơn giản hơn.
\end{itemize}
\end{phuongphap}

\subsection{Underfitting}

\begin{dinhnghia}[title={Underfitting (Thiếu khớp)}]
Underfitting xảy ra khi mô hình quá đơn giản và không thể nắm bắt pattern/mối quan hệ trong dữ liệu.
Điều này dẫn tới hiệu năng kém trên cả train và test.
\end{dinhnghia}

\begin{phuongphap}[title={Cách giảm underfitting}]
\begin{itemize}
  \item Tăng độ phức tạp mô hình.
  \item Thu thập thêm dữ liệu.
  \item Giảm regularization.
  \item Feature engineering.
\end{itemize}
\end{phuongphap}

\begin{luuy}[title={Cân bằng giữa underfitting và overfitting}]
Giảm underfitting là một bài toán cân bằng giữa độ phức tạp mô hình và lượng dữ liệu hiện có.
Tăng độ phức tạp có thể giúp bớt underfitting, nhưng nếu dữ liệu không đủ để “đỡ” độ phức tạp đó, mô hình có thể chuyển sang overfitting.
Vì vậy cần theo dõi hiệu năng và điều chỉnh độ phức tạp khi cần.
\end{luuy}

\subsection{Khi nào nên “làm cho máy học”?}

\begin{dongco}[title={Vì sao phải để máy học thay vì viết quy tắc?}]
Ngoài lý do “cần ML”, còn có câu hỏi: trong kịch bản nào ta nên làm cho máy học?
Có nhiều tình huống cần ra quyết định dựa trên dữ liệu với hiệu quả và quy mô lớn.
\end{dongco}

\begin{phuongphap}[title={Một số tình huống}]
\begin{itemize}
  \item \textbf{Thiếu chuyên gia con người}: ví dụ điều hướng ở vùng chưa biết hoặc nhiệm vụ trong không gian.
  \item \textbf{Tình huống động}: hành vi thay đổi theo thời gian; ví dụ kết nối mạng, hạ tầng tổ chức.
  \item \textbf{Khó chuyển chuyên môn thành bài toán tính toán}: con người có trực giác/chuyên môn nhưng không viết được quy tắc; ví dụ nhận dạng giọng nói, tác vụ nhận thức.
\end{itemize}
\end{phuongphap}

\subsection{Định nghĩa học máy của Tom Mitchell và mô hình T--P--E}

\begin{dinhnghia}[title={Định nghĩa (Tom Mitchell)}]
Một chương trình máy tính được gọi là “học” từ kinh nghiệm \(E\) đối với một lớp nhiệm vụ \(T\) và thước đo hiệu năng \(P\),
nếu hiệu năng của nó khi thực hiện nhiệm vụ trong \(T\), được đo bằng \(P\), được cải thiện nhờ kinh nghiệm \(E\).
\end{dinhnghia}

\begin{ynghia}[title={Ba thành phần của bài toán học}]
Định nghĩa trên nhấn mạnh ba thành phần (cũng là thành phần chính của thuật toán học):
\begin{itemize}
  \item \textbf{Task (T)}: nhiệm vụ.
  \item \textbf{Performance (P)}: thước đo hiệu năng.
  \item \textbf{Experience (E)}: kinh nghiệm (dữ liệu/quá trình học).
\end{itemize}
\end{ynghia}

\begin{phuongphap}[title={Cách diễn giải ngắn}]
ML là lĩnh vực AI gồm các thuật toán học mà:
\begin{itemize}
  \item Cải thiện hiệu năng \(P\).
  \item Khi thực hiện nhiệm vụ \(T\).
  \item Theo thời gian nhờ kinh nghiệm \(E\).
\end{itemize}
\end{phuongphap}

\subsubsection{Task (T)}
\begin{ynghia}[title={Task theo góc nhìn ML}]
Theo góc nhìn “bài toán thực”, \(T\) có thể là dự đoán giá nhà hoặc chọn chiến lược marketing tốt.
Theo góc nhìn ML, \(T\) thường là những nhiệm vụ mà cách lập trình truyền thống khó giải vì phải xử lý trên nhiều điểm dữ liệu.
Ví dụ về nhiệm vụ ML: classification, regression, structured annotation, clustering, transcription, \ldots
\end{ynghia}

\subsubsection{Experience (E)}
\begin{ynghia}[title={Kinh nghiệm là gì?}]
Kinh nghiệm là kiến thức thu được từ các điểm dữ liệu cung cấp cho thuật toán.
Khi được cung cấp dataset, mô hình chạy lặp và học các pattern nội tại; phần học được gọi là experience \(E\).
Theo tương tự với con người: con người học từ tình huống, quan hệ, \ldots
Supervised/unsupervised/reinforcement là các cách khác nhau để “tạo kinh nghiệm”.
Kinh nghiệm của mô hình sẽ được dùng để giải nhiệm vụ \(T\).
\end{ynghia}

\subsubsection{Performance (P)}
\begin{ynghia}[title={Thước đo hiệu năng}]
Thước đo hiệu năng cho biết thuật toán có đang làm đúng kỳ vọng không.
\(P\) là metric định lượng đánh giá mô hình thực hiện \(T\) bằng kinh nghiệm \(E\) tốt đến đâu.
Một số metric: accuracy score, F1 score, confusion matrix, precision, recall, sensitivity, \ldots
\end{ynghia}

